\subsection{Flasqueness and the cosimplicial obstruction}

In this section the coefficient ring $k$ is constant. We distinguish
flasqueness in the sense of SGA~4, V, Definition~4.1, from absolute
$\Gamma$-acyclicity. The former means
$H^q(\widehat\Delta/X,A|_X)=0$ for every simplicial set $X$ and $q>0$.
We say \emph{tensor-universally flasque} if $A\otimes_k M$ is flasque
for every simplicial $k$-module $M$. The condition on cote~136,
scan~80, explicitly includes tensoring with every module sheaf.

\begin{theorem}[C01: \status{Proved}; L6]\label{sa08:C01}
Let $k\ne0$ be a commutative ring. In the topos of cosimplicial sets
there is no monomorphism $k\to\Phi$ such that $\Phi\otimes_k M$ is
$\Gamma$-acyclic for every cosimplicial $k$-module $M$. In particular,
the conditions of cote~136, scan~80, cannot hold there, confirming the
\emph{Non} on scan~91. Flatness of the cokernel and the conditions on
higher exterior and divided powers are unnecessary for this obstruction.
\end{theorem}
\begin{proof}
Let $V$ be the cosimplicial denormalization of the cochain complex
$k[1]$. Thus $V^0=0$, $V^1=k$, and its generator $v$ is a normalized
one-cocycle. Write $t_n$ for the image of $1$ in $\Phi^n$.
Naturality of the monomorphism gives $d^i t_1=t_2$ for all three
cofaces into degree two. Consequently
\[
 d(t_1\otimes v)=t_2\otimes(d^0v-d^1v+d^2v)=0.
\]
This is a nonzero cocycle because $\Phi^1\otimes_k V^1=\Phi^1$ and
$t_1\ne0$. There are no degree-zero cochains in $\Phi\otimes_k V$,
so it is not a boundary. Normalized cosimplicial cohomology computes
derived global sections: under the dual Dold--Kan equivalence,
global sections are $\ker(d:C^0\to C^1)$, and the projective disk
resolution of $k[0]$ computes its right derived functors as $H^q(C)$.
Thus $H^1(\Phi\otimes_k V)\ne0$.
\end{proof}




\begin{theorem}[C02: \status{Proved}; exact replacement for L5]\label{sa08:C02}
For a simplicial $k$-module $A$, flasqueness is equivalent to
acyclicity of the homological normalized complex $N(A)$, including
$H_0(N(A))=0$. If every $A_n$ is projective over $k$, these conditions
imply an additive simplicial contraction. Without projectivity this
last implication is false, even if all terms are flat.
\end{theorem}
\begin{proof}
For the projection $p_X:\widehat\Delta/X\to\widehat\Delta$, the
left adjoint $p_{X!}$ on module sheaves is pointwise direct sum over
the fibers, hence exact. Thus $p_X^*$ preserves injectives, and
adjunction gives
\[
 H^q(\widehat\Delta/X,p_X^*A)
 =\Ext^q_{s\operatorname{Mod}_k}(k[X],A)
 =\Ext^q_{\operatorname{Ch}_{\ge0}(k)}(N(k[X]),N(A)).
\]
The last equality uses the exact Dold--Kan equivalence
\cite[Theorem~14.24.3]{stacks-doldkan}. The complex $P=N(k[X])$ is free
in every degree, on the nondegenerate simplices of $X$.

For a degreewise projective complex $P$, let $D^n(P_n)$ be the identity
disk in homological degrees $n,n-1$. The projective epimorphism
\[
 D(P)=P_0[0]\oplus\bigoplus_{n\ge1}D^n(P_n)\longrightarrow P
\]
is $(a,b)\mapsto a+d_Pb$ in degree $n$, where
$D(P)_n=P_n\oplus P_{n+1}$. Its kernel is $P_{n+1}$ in degree $n$,
with differential $-d_P$. Iteration produces a projective resolution.
Applying $\Hom(-,N(A))$ identifies its positive Ext groups with the
positive cohomology of
\[
 \underline\Hom(P,N(A))^q
 =\prod_{n\ge0}\Hom_k(P_{n+q},N(A)_n).
\]
If $N(A)$ is exact, a cocycle $f$ in degree $q>0$ is a boundary.
To see this directly, set the degree $q-1$ initial homotopy component
to zero and choose successively $g_j:P_j\to N(A)_{j-q+1}$ satisfying
\[
 d_Ag_j=f_j+(-1)^{q-1}g_{j-1}d_P.
\]
The right side is a cycle by the cocycle equation and the previous
identity. Exactness and projectivity of $P_j$ give the lift. This
recursive construction uses no finiteness of $X$.

Conversely, for $X=\Delta[m+1]/\partial\Delta[m+1]$, one has
$N(k[X])=k[0]\oplus k[m+1]$. The displayed resolution gives
\[
 H^1(\widehat\Delta/X,p_X^*A)=H_m(N(A)),\qquad m\ge0.
\]
Thus flasqueness implies exactness. If the terms of $N(A)$ are
projective, its short exact cycle sequences split inductively from
degree zero: $Z_0=N(A)_0$ is projective, the surjection onto $Z_0$
splits, its kernel $Z_1$ is projective, and so on. These splittings
give a chain contraction. Dold--Kan transports it to an additive
simplicial homotopy. Each normalized term is a direct summand of
$A_n$, so the stated projectivity hypothesis suffices. The failures
without projectivity are proved next.
\end{proof}

\begin{proposition}[C03: \status{Refuted}; L5 as an additive-contraction assertion]\label{sa08:C03}
An object acyclic on every slice need not admit an additive simplicial
contraction. This fails for a finite normalized complex.
\end{proposition}
\begin{proof}
Take $k=\Z$ and the exact complex
\[
 0\longrightarrow\Z/2\xrightarrow{\;2\;}\Z/4
 \xrightarrow{\bmod2}\Z/2\longrightarrow0
\]
in homological degrees $2,1,0$. Its Dold--Kan denormalization is
flasque by Statement~\ref{sa08:C02}. A contraction would give a
section $\Z/2\to\Z/4$ of reduction modulo two. The image of $1$
under any such group homomorphism is $0$ or $2$, both reducing to
zero, so no section exists. The program
\texttt{code/C\_finite\_obstructions.py} checks exactness and exhausts
these two possible homomorphisms; its complete output is supplied.
The underlying simplicial set is weakly contractible. The false
conclusion here is the \emph{additive} contraction needed to apply
nonadditive module functors to a homotopy.
\end{proof}

\begin{proposition}[C04: \status{Proved}; counterexample to the flat reading of scan~100]\label{sa08:C04}
There is a tensor-universally flasque simplicial abelian group whose
terms are flat over $\Z$ but which admits no additive contraction.
\end{proposition}
\begin{proof}
Let $b_j$ be $1$ when $j$ is a perfect square and $0$ otherwise, and put
$a_n=\sum_{0\le j<n}b_j2^j$. In $\Q^2$ put
\[
 z=(1,0),\qquad x_n=(a_n/2^n,1/2^n),\qquad
 E=\langle z,x_0,x_1,\ldots\rangle_{\Z}.
\]
The relations are $2x_{n+1}=x_n+b_nz$. Projection to the second
coordinate gives an exact sequence
\[
 0\longrightarrow\Z z\longrightarrow E\longrightarrow\Z[1/2]
 \longrightarrow0.
\]
Indeed a finite expression can be rewritten as $cz+mx_N$, and its
second coordinate vanishes only if $m=0$. All three terms are
torsionfree, hence flat over $\Z$.

A section must send $1$ to $x_0+cz$ for an integer $c$. For its
$2^n$-th part to belong to $E$, one must have
$c\equiv a_n\pmod{2^n}$. To check the necessity without assuming a
presentation, rewrite any putative part using $x_N$ with $N\ge n$;
then $a_N\equiv a_n\pmod{2^n}$ gives the same congruence. These
congruences are impossible: the binary digits $b_j$ are neither
eventually zero nor eventually one, whereas the $2$-adic expansion
of an ordinary integer has one of these two forms.

Place the sequence in degrees $2,1,0$ and denormalize. It is flasque
by Statement~\ref{sa08:C02}; a contraction would split the sequence.
Its terms and cycles are flat, so tensoring its normalized complex
with any nonnegative complex is exact, by the first-quadrant tensor
filtration. Eilenberg--Zilber transfers this to the levelwise
simplicial tensor product. It is therefore tensor-universally flasque.
\end{proof}

The supplied check event S3-136-002-141-CE036 uses the constant
simplicial module $k$ as a counterexample to an ``absolute flasqueness''
reading. Its calculation of absolute $\Gamma$-acyclicity is correct,
but does not refute the SGA~4 meaning of flasqueness. Indeed global
sections in $\widehat\Delta$ are evaluation at degree zero, hence
exact for \emph{every} simplicial module. The constant module is not
flasque: Statement~\ref{sa08:C02}, with $m=0$, detects it on the
simplicial circle. Statements~\ref{sa08:C02} and \ref{sa08:C04} handle
both the ``projective'' and ``flat'' readings of the overwritten word
on scan~100.

\begin{theorem}[C05: \status{Proved}; repair of the conclusion on scan~100]\label{sa08:C05}
Let $A$ be a degreewise flat simplicial $k$-module with $N(A)$ acyclic.
Then for every $r>0$, both $\Gamma_k^r(A)$ and $\Lambda_k^r(A)$ are
flat and tensor-universally flasque.
\end{theorem}
\begin{proof}
The cycles of $N(A)$ are flat, inductively from $Z_0=N(A)_0$, using
$0\to Z_n\to N(A)_n\to Z_{n-1}\to0$. Christensen--Holm's theorem
\cite[Theorem~7.3 in the opened preprint]{christensen-holm}
expresses such a complex as a filtered colimit of bounded
contractible complexes of finite free modules. To work in
nonnegative degrees, apply the smart truncation that replaces the
degree-zero term by its kernel into degree minus one and discards
negative degrees. This truncation commutes with filtered colimits.
On a split contractible complex it remains contractible and consists
of finite projective modules. Dold--Kan commutes with filtered
colimits.

An ordinary functor preserves simplicial homotopy identities, which
are identities of compositions; a reduced functor also preserves
the zero map. Positive divided and exterior powers are reduced,
commute with filtered colimits, and take finite free modules to
finite free modules. They take projectives to projectives because
they preserve retracts. Consequently their values on each
contractible approximant are simplicially contractible and
levelwise flat. Exactness of filtered colimits proves acyclicity
and flatness of their values on $A$.

Finally, a nonnegative exact complex of flat modules has flat cycles
by the same induction. Tensoring with an arbitrary nonnegative
complex is exact, using the first-quadrant filtration. The
Eilenberg--Zilber chain equivalence, together with
Statement~\ref{sa08:C02}, proves flasqueness after every tensor product.
\end{proof}

\subsection{The special category and the uniqueness conjecture}

\begin{theorem}[C06: \status{Proved}; repaired scan~98 proposition]\label{sa08:C06}
For $\widehat\Delta$ and any constant commutative ring $k$, the
category of special tensor-universally flasque PD de Rham structures
of scan~94 has contractible nerve. Its localization at all morphisms
is equivalent to the terminal category. The standard elementary
theory maps to every object, so it is an admissible common source
for every pair of objects in this special category.
\end{theorem}
\begin{proof}
For a special structure $A$, let $\Phi_A=A^{1,0}$ and
$\Psi_A=A^{1,1}$. Specialness gives the pure exact sequence
$0\to kT\to\Phi_A\to\Psi_A\to0$ with flat terms. The elementary
structure
\[
 E(A)=\Gamma_k^*(\Phi_A)\otimes_k\Lambda_k^*(\Psi_A)
\]
with its standard differential maps naturally to $A$ by the free
PD-alternating universal property. We recall why it resolves
$S=k\{T\}$. Express the flat quotient $\Psi_A$ as a filtered
colimit of finite free modules and pull back the extension. Each
pulled-back extension splits, since its quotient is projective.
For a split extension, the resulting complex is $S$ tensored with
divided-power Poincare complexes. In one variable the latter has
$d(x^{(n)})=x^{(n-1)}dx$ and contraction
$x^{(n-1)}dx\mapsto x^{(n)}$; tensoring these free complexes proves
the assertion in finitely many variables. Filtered colimits give the
general assertion. Statement~\ref{sa08:C05} makes every term of
positive complementary degree universally flasque.

Write $\Phi_{\mathrm{std},n}=k^{n+1}$. The right adjoint of
evaluation at zero is $i_{0*}(M)_n=M^{n+1}$. Its adjunction unit gives
$\Phi_A\to i_{0*}((\Phi_A)_0)$; explicitly, it evaluates a simplex
on its vertices. Define
\[
 Q(A)=\mathrm{DR}\bigl(i_{0*}((\Phi_A)_0),T\bigr).
\]
The quotient of the marked copy of $k$ in degree $n$ is a direct sum
of $(\Phi_A)_0/kT$ and $n$ copies of $(\Phi_A)_0$, hence flat.
The coinduced module is simplicially contractible and flat. Thus
$Q(A)$ belongs to the same category. There are natural morphisms
\[
 A\longleftarrow E(A)\longrightarrow Q(A)
 \longleftarrow\mathrm{DR}(\Phi_{\mathrm{std}},T).
\]
After localization they give a natural isomorphism from the
constant functor at the last object to the identity. If
$\eta_A$ denotes that isomorphism, naturality forces every arrow
$f:A\to B$ to equal $\eta_B\eta_A^{-1}$. All Hom sets therefore
contain exactly one arrow. Before localization the same natural
transformations give a contraction of the nerve.

For the final assertion, $N(\Phi_{\mathrm{std}})$ is the identity
disk $D^1(k)$ with its marked element in degree zero. Since
$\Phi_A$ is simplicially acyclic, lift its marked element
$T\in(\Phi_A)_0$ through $N_1(\Phi_A)\to(\Phi_A)_0$.
The lift defines a marked simplicial map
$\Phi_{\mathrm{std}}\to\Phi_A$, hence a morphism
$\mathrm{DR}(\Phi_{\mathrm{std}},T)\to E(A)\to A$.
Doing this for each of two objects proves the common-source claim.
\end{proof}

For a general ringed topos, the corresponding conditional statement
requires nonemptiness in addition to preservation of flat universal
flasqueness by positive divided powers. Under those hypotheses one
can use a fixed admissible embedding $k\to\Phi_0$ and the common
target $(\Phi\amalg_k\Phi_0)\otimes_k\Phi_0$ in place of the
coinduced construction. The two natural maps give the same
contraction. Without nonemptiness the category is empty, not
terminal; Statement~\ref{sa08:C01} explains why that distinction
matters.

\paragraph{The category used for C.3.}
We make the explicit conditions of scans~111, 117 and~119
self-contained. Put $S=\Gamma_k(T)$. An object is a simplicial
bigraded, graded-commutative alternating PD differential
$S$-algebra $A_n^{w,q}$, with $w,q\geq0$, differential of
bidegree $(0,1)$, the specified coefficient map
$S\to A_n^{*,0}$, and the compatible $k$-augmentation.
For each $d\geq0$ there must be bounds $\lambda(d),\mu(d)$
such that the following two conditions hold:
\begin{enumerate}
\item for every simplicial degree $n$ and every $w\geq\lambda(d)$,
the coefficient map $S^w\to H^0(A_n^{w,*})$ is an isomorphism
and $H^q(A_n^{w,*})=0$ for $1\leq q\leq d$;
\item for every $w\geq\mu(d)$ and $0\leq q\leq d$, the
homological normalized complex $N(A_*^{w,q})$ is exact,
including degree zero.
\end{enumerate}
Morphisms preserve both gradings, differential, products,
divided powers, coefficients and augmentation. By
Statement~\ref{sa08:C02}, the second condition is precisely
eventual slice flasqueness. Bounds are uniform in $n$ in the
first condition. We call these objects \emph{admissible} in
the C.3 statements below. When we require termwise $k$-flatness,
or the support condition $q\leq w$, we state that additional
restriction explicitly.

\begin{theorem}[C07: \status{Refuted}; strict equalization, scan~121(b)]\label{sa08:C07}
For every nonzero commutative ring $k$, the claim that every parallel
pair between admissible generalized PD de Rham theories can be
equalized by an admissible common source is false. It fails even
when both theories are special and tensor-universally flasque.
\end{theorem}
\begin{proof}
Let $S=k\{T\}$ and $F=\mathrm{DR}(\Phi,T)$ be the standard theory,
where $\Phi_n=k^{n+1}$. Put
\[
 W=\Phi\amalg_k\Phi,\qquad V=\Phi\otimes_k W,
 \qquad B=\mathrm{DR}(V,T\otimes T).
\]
The map $j:W\to V$, $w\mapsto T\otimes w$, is split in each
simplicial degree, using evaluation at any chosen vertex in the first
factor. All its terms and cokernels are free. The module $V$ is
acyclic: the first factor $\Phi$ is contractible and flat, so
Eilenberg--Zilber, or the tensor of its simplicial contraction with
the constant identity homotopy of $W$, proves the claim. Also
$V/k(T\otimes T)$ is flat, since the marked vector is unimodular
in every degree. Statement~\ref{sa08:C05} and the relative
Poincare argument in Statement~\ref{sa08:C06} show that $B$ is
special and tensor-universally flasque.

The map $j$ induces a monomorphism
\[
 F\otimes_S F=\mathrm{DR}(W,T)\lhook\joinrel\longrightarrow B.
\]
To verify injectivity over arbitrary $k$, split $j$ degreewise and
split off its common marked line. Divided powers and exterior
powers of the resulting inclusions into direct sums are injective,
by their direct-sum decompositions. Let $i_1,i_2:F\to B$ be the
two inclusions through $F\otimes_S F$.

In each simplicial degree, $F_n$ is free over $S$ on divided-power
and exterior monomials, including $1$. Comparison of the tensor
basis proves that $a\otimes1=1\otimes a$ holds exactly for $a\in S$.
The injection above therefore implies
\[
 \operatorname{Eq}(i_1,i_2)=S
\]
in simplicial bigraded differential algebras. If $h:H\to F$ from
an admissible theory equalized the pair, it would factor through the
constant coefficient algebra $S$. Its map on positive exterior-degree
cohomology after sections over any $X$ would consequently be zero.
But a morphism of admissible theories preserves the augmentation
and the two comparison arrows of scans~113--115, and hence induces
the ordinary-cohomology isomorphism after the shadow functor.
Taking $X=\Delta[1]/\partial\Delta[1]$, whose degree-one cohomology
is $k\ne0$, gives a contradiction.
\end{proof}

This counterexample does not refute connectedness or simple
connectedness of the localized category. A calculus of common-source
fractions would imply that property, but is not equivalent to it.
The claimed equivalence implicit in scan~121 must therefore be
separated from the groupoid question.

\begin{theorem}[C08: \status{Proved under stated hypotheses}; rational uniqueness]\label{sa08:C08}
Let $k$ be a $\Q$-algebra. Consider the category of flat bigraded
PD de Rham theories satisfying the uniform eventual resolution
condition of scan~117 and the eventual flasqueness condition of
scan~119. Its nerve is contractible, and its localization at all
morphisms is terminal.
\end{theorem}
\begin{proof}
Let $F$ be the standard theory. We prove that $F\otimes_S C$ remains
in the category for every $C$ in it. Filter $F$, as a graded
$S$-module, by polynomial degree in $\Phi_{\mathrm{std}}/kT=\Psi$.
In each simplicial degree this filtration is split over $S$, so
it remains exact on tensoring with $C$. Its associated graded in
weight $p$ and exterior degree $q$ has terms
\[
 C^{p-r-s,q-s}\otimes_k\operatorname{Sym}^r(\Psi)
 \otimes_k\Lambda^s(\Psi),
 \qquad r,s\ge0,\quad s\le q.
\]
The filtration in a fixed weight is finite. For $r\ge2$ the symmetric
factor is acyclic and flat: $N(\Psi)=k[1]$, and the same averaging and
sign argument as above works over any $\Q$-algebra. It is therefore
tensor-universally flasque by Statement~\ref{sa08:C05} and its proof.
For $r=0,1$, all $C$-terms involved have bounded exterior degree and
weight tending to infinity with $p$; eventual flasqueness of $C$
therefore applies. The remaining factors are flat. Finite
devissage proves eventual flasqueness of $F\otimes_S C$.

Simplexwise the polynomial de Rham complex $F_n$ has an $S$-linear
contraction to $S$, preserving total weight. Choose vertex
coordinates and use the radial contraction, dividing by the
positive sum of polynomial and exterior degrees; the denominators
are units in $k$. Tensoring this contraction with $C_n$ shows
$C_n\to F_n\otimes_S C_n$ is a weightwise homotopy equivalence.
The uniform eventual resolution bounds are consequently unchanged.
All terms remain flat by the displayed split filtration.
Finally, the natural zigzag
\[
 C\longrightarrow F\otimes_S C\longleftarrow F
\]
contracts the nerve and the localization exactly as in
Statement~\ref{sa08:C06}.
\end{proof}

Outside the rational case, the associated graded uses
$\Gamma^r(\Psi)$ instead of $\operatorname{Sym}^r(\Psi)$;
these divided-power factors need not be acyclic.
The next argument uses different common sources and proves the
localized statement over a wider class of coefficient rings.

\begin{theorem}[C09: \status{Proved under stated hypotheses}; uniqueness over a localized coefficient ring]\label{sa08:C09}
Fix a prime $p$, and let $k$ be a $\Z_{(p)}$-algebra. In the category
of generalized PD de Rham theories satisfying the explicit eventual
conditions of scans~117 and~119, the localization at all morphisms
is terminal. Any two objects have an admissible common source.
No assertion of strict equalization is made; that assertion is
false by Statement~\ref{sa08:C07}.
\end{theorem}
\begin{proof}
Write $S=k\{T\}$. For $N=p^e$ let $R_N=k\{U\}$, assigning $U$
weight $N$, and map $R_N\to S$ by $U\mapsto T^{(N)}$. The graded
$R_N$-module $S$ is free on $T^{(r)}$, $0\le r<N$. Indeed
\[
 U^{(m)}T^{(r)}\longmapsto
 \frac{(mN)!}{m!(N!)^m}\binom{mN+r}{r}T^{(mN+r)}.
\]
Both coefficients are prime to $p$. For the first, Legendre's
formula and $s_p(N)=1$, $s_p(mN)=s_p(m)$ give valuation zero; for
the second, the base-$p$ digits of $mN$ and $r<N$ do not overlap.
Thus every displayed coefficient is a unit in $k$, and each weight
has exactly one displayed basis vector.

Let $\Phi_n=k^{n+1}$ with marked constant vector $1$, and define
\[
 G_N=S\otimes_{R_N}\mathrm{DR}(\Phi,U),
\]
where the generators of $\Phi$ and of its quotient $\Psi$ have
weight $N$. This is an admissible theory. Simplexwise it resolves
$S$, because $S$ is free over $R_N$ and the relative PD complex
resolves $R_N$. As a graded simplicial module it is the finite direct
sum of the usual elementary complex, with weights multiplied by
$N$, shifted by $r=0,\ldots,N-1$. In weight $n$ and exterior degree
$q$, it is therefore flasque whenever
$q<\lfloor n/N\rfloor$. This verifies the eventual condition
uniformly.

The normalized complex of $\Phi$ is the identity disk $D^1(k)$,
so $\Phi$ is projective in simplicial $k$-modules. If $A$ is
admissible and $A^{N,0}$ is flasque, its normalized complex is
acyclic by Statement~\ref{sa08:C02}. Lift the marked element
$T^{(N)}\in A_0^{N,0}$ through
$d:N_1(A^{N,0})\to A_0^{N,0}$. The lift gives a simplicial linear
map $\Phi\to A^{N,0}$ carrying the marked vector to $T^{(N)}$.
The free PD-alternating universal property, followed by base change,
gives a morphism $G_N\to A$. The eventual condition ensures this
works for all sufficiently large powers $N$ of $p$. It proves the
common-source assertion.

We next show that any two maps $f,g:G_N\to A$ become equal after
localization, assuming $A^{N,0}$ is flasque. Their restrictions to
$\Phi$ agree on the marked line, hence give a linear map
$W=\Phi\amalg_k\Phi\to A^{N,0}$. Embed $W$ in
$V=\Phi\otimes_k W$ by $w\mapsto1\otimes w$. Its cokernel is
free in every simplicial degree, so its normalized terms are
projective. The Ext computation in the proof of
Statement~\ref{sa08:C02} shows
\[
 \Ext^1_{s\operatorname{Mod}_k}(V/W,A^{N,0})=0.
\]
Thus the map on $W$ extends to $V$. By the free universal property
it yields a map from
\[
 H_N=S\otimes_{R_N}\mathrm{DR}(V,U)
\]
to $A$. This theory is admissible by the same finite free
base-change argument, since $V$ is flat and acyclic.

The fold map $W\to\Phi$ sends the two copies identically to
$\Phi$. Compose $\Phi\otimes W\to\Phi\otimes\Phi$ with
pointwise multiplication of vertex functions. It is a simplicial
linear retraction of both inclusions $\Phi\to V$, and induces a
common retraction $H_N\to G_N$ of the two inclusions
$G_N\rightrightarrows H_N$. Once this retraction is inverted, the
two inclusions coincide. Their composites to $A$ are $f$ and $g$,
so $f=g$ in the localization.

Fix the standard theory $F$. For each $A$ choose a sufficiently
large $N$ and maps $a_N:G_N\to F$, $f_N:G_N\to A$, and set
$\eta_A=f_Na_N^{-1}$ in the localization. For two choices $N,N'$,
choose a larger power $L$ for which all needed degree-zero terms
are flasque. There are maps $G_L\to G_N$ and $G_L\to G_{N'}$ by
the preceding lifting argument. After composition, the two maps
to $A$, and likewise the two maps to $F$, coincide in the
localization by the previous paragraph. Hence $\eta_A$ is
independent of every choice. For any morphism $u:A\to B$, use a
single sufficiently large $N$; the same argument gives
$u\eta_A=\eta_B$. Thus $\eta$ is a natural isomorphism from the
constant functor at $F$ to the identity. Every morphism in the
localized groupoid equals $\eta_B\eta_A^{-1}$, proving terminality.
\end{proof}

The remaining unrestricted coefficient-ring gate is now narrower:
the argument proves uniqueness over every $\Z_{(p)}$-algebra,
including every field and every $\Q$-algebra, but does not descend
these local comparisons to an arbitrary ring such as $\Z$. A
single weight multiplier $N>1$ is not a power of every residual
prime, so the finite-free comparison used above cannot be applied
globally without another construction.

\begin{corollary}[C10: \status{Proved under stated hypotheses}]\label{sa08:C10}
Statement~\ref{sa08:C09} holds for every ring of positive
characteristic, with the same explicit eventual-condition category.
\end{corollary}
\begin{proof}
If $\operatorname{char}(k)=m>0$, write
$m=\prod_{p\mid m}p^{a_p}$. The Chinese remainder idempotents in
$\Z/m$ decompose $k$ as a finite product $\prod_{p\mid m}k_p$,
where each $k_p$ is a $\Z_{(p)}$-algebra. Every graded module,
differential algebra, divided-power structure and morphism over $k$
splits under these central degree-zero idempotents. The two eventual
conditions hold over $k$ exactly when they hold in each factor,
because there are finitely many factors and one takes maxima of
their bounds. Thus the category is the finite product of the
corresponding categories over $k_p$. Its groupoid localization is
the product of their localizations: pad finite zigzags with identity
arrows to represent every componentwise tuple, and use the universal
property for the converse. Each factor is terminal by
Statement~\ref{sa08:C09}. Common sources are also taken factorwise.
\end{proof}

\begin{proposition}[C11: \status{Refuted}; tensor closure in the integral uniqueness argument]\label{sa08:C11}
The tensor product over $S=k\{T\}$ of two admissible special
theories need not satisfy eventual flasqueness. This fails for
$F\otimes_S F$, where $F$ is the standard theory, both over $\Z$
and over every $\F_p$. Thus the tensor construction used in
Statement~\ref{sa08:C08} cannot be applied over arbitrary rings.
\end{proposition}
\begin{proof}
The degree-zero simplicial module $F^0=\Gamma_k(\Phi)$ is free over
$S$ in every simplicial degree: splitting the marked line in
$\Phi_n$ gives $F_n^0=S\otimes_k\Gamma_k(k^n)$. Its normalized
terms are projective over $S$, as direct summands of free modules.
By Statement~\ref{sa08:C05}, its positive-weight homotopy vanishes,
while its weight-zero homotopy is $k$ in degree zero. Thus it is a
projective $S$-resolution of $k=S/S^+$. Eilenberg--Zilber over $S$
gives
\[
 \pi_i(F^0\otimes_S F^0)=\operatorname{Tor}_i^S(k,k),\qquad
 \operatorname{Tor}_1^S(k,k)=S^+/(S^+)^2.
\]
The second equality follows by tensoring
$0\to S^+\to S\to k\to0$ with $k$.

In weight $p^e$ over $\F_p$, all the internal binomial coefficients
$\binom{p^e}{a}$ vanish, so $T^{(p^e)}$ gives a nonzero element of
$S^+/(S^+)^2$. Such weights are unbounded; eventual flasqueness
fails in exterior degree zero.

Over $\Z$, its weight-$n$ component is $\Z/d_n$, where
$d_n=\gcd_{0<a<n}\binom na$ for $n>1$. A prime $p$ divides every
internal coefficient exactly when
$(1+x)^n=1+x^n$ in $\F_p[x]$, equivalently when $n$ is a power
of $p$, by expansion in base $p$. If $n=p^e$, then
\[
 \binom{p^e}{p^{e-1}}
 =p\binom{p^e-1}{p^{e-1}-1}
\]
has $p$-adic valuation one: all base-$p$ digits of $p^e-1$ are
$p-1$, so the second factor is a unit modulo $p$. Thus
$d_{p^e}=p$. The resulting unbounded nonzero $\Z/p$ components
prove the integral failure as well.
\end{proof}

\begin{proposition}[C12: \status{Proved}; obstruction to arithmetic gluing of the local sources]\label{sa08:C12}
For each prime $p$ choose a power $N_p=p^{e_p}$ and let
$G_{N_p}$ be the theory of Statement~\ref{sa08:C09} over
$\Z_{(p)}$. Suppose that the integers $N_p$ are unbounded.
There is no integral admissible theory $H$ with coefficient-preserving
morphisms
\[
 H\otimes_{\Z}\Z_{(p)}\longrightarrow G_{N_p}
 \qquad\text{for every prime }p.
\]
Consequently, arithmetic pullbacks, including homotopy pullbacks,
that retain projections to all these local sources cannot be
admissible. Changing the rational path object cannot remove this
obstruction while retaining those projections.
\end{proposition}
\begin{proof}
Choose $W$ so that $H^{w,0}$ is simplicially acyclic for $w\ge W$.
For any such $w$, the marked element $T^{(w)}\in H_0^{w,0}$ is a
normalized boundary. Localization at $p$ is exact, so it remains a
boundary after tensoring with $\Z_{(p)}$.

Choose $p$ with $N_p>w$. The finite-free description in
Statement~\ref{sa08:C09} gives
\[
 G_{N_p}^{w,0}=\Z_{(p)}T^{(w)}
 \qquad(0\le w<N_p)
\]
as a constant simplicial module. Its marked element is nonzero in
$\pi_0$. A coefficient-preserving simplicial map cannot send the
boundary $T^{(w)}$ to this nonboundary. This proves the assertion.
The argument only uses the putative projection maps, so applies to
every proposed pullback with those projections.
\end{proof}
